% Generated draft: requires mathematical and source review.
\documentclass[11pt]{article}
\usepackage[margin=25mm]{geometry}
\usepackage{fontspec}
\setmainfont{Latin Modern Roman}
\usepackage{amsmath,amssymb,mathtools,unicode-math}
\setmathfont{Latin Modern Math}
\usepackage{enumitem,needspace,xurl,hyperref}
\setlength{\parindent}{0pt}
\setlength{\parskip}{5pt}
\setlength{\emergencystretch}{3em}
\newcommand{\N}{\mathbb{N}}
\newcommand{\Z}{\mathbb{Z}}
\newcommand{\Q}{\mathbb{Q}}
\newcommand{\R}{\mathbb{R}}
\newcommand{\C}{\mathbb{C}}
\newcommand{\F}{\mathbb{F}}
\newcommand{\A}{\mathbb{A}}
\newcommand{\PP}{\mathbb{P}}
\newcommand{\E}{\mathbb{E}}
\newcommand{\eps}{\varepsilon}
\newcommand{\dd}{\,\mathrm{d}}
\begin{document}
\section*{UnsolvedMath ID: 30000032}
\textbf{Source problem code:} catalogue entry
\subsection{Definitions and mathematical statement}
The verified clean statement contains two questions:
\begin{enumerate}
\item Describe the nonlinear spectral behavior near an eigenvalue $\lambda_0$ that is not simple for the operator considered in the report.
\item Characterize a class of homogeneous maps on $\R^N$ for which every eigenvalue is isolated.
\end{enumerate}
The original short extraction retained only the first question; the supplied verification says the second was restored from the same open-problem paragraph.

For clarity, in the usual degree-one interpretation a homogeneous map is $F:\R^N\to\R^N$ satisfying $F(tx)=tF(x)$ for the prescribed positive scalars $t>0$. Its real eigenvalue set is
\[
 \sigma_{\mathrm{eig}}(F)=\{\lambda\in\R:\exists x\in\R^N\setminus\{0\},\ F(x)=\lambda x\}.
\]
An eigenvalue $\lambda$ is isolated if some $\varepsilon>0$ satisfies
\[
 \sigma_{\mathrm{eig}}(F)\cap(\lambda-\varepsilon,\lambda+\varepsilon)=\{\lambda\}.
\]
Thus the second target seeks structural hypotheses on a class $\mathcal H$ implying
\[
 \forall F\in\mathcal H\ \forall\lambda\in\sigma_{\mathrm{eig}}(F),
 \qquad\lambda\text{ is isolated in }\sigma_{\mathrm{eig}}(F).
\]
This concerns isolation of spectral values, not uniqueness or isolation of eigenvectors. The degree-one positive-homogeneity convention is made explicit for interpretation; the extract does not specify the degree, negative-scalar behavior, regularity, or other intended restrictions on the class.

For the first question, the relevant operator, its domain, a parameter family, and the meaning of nonlinear eigenvalue multiplicity are absent. Consequently no particular Fu\v cik operator, perturbation equation or branching theorem can be supplied without inventing the missing setup. ``Not simple'' must be matched to the source's multiplicity notion rather than defined by a linear-eigenspace formula for an arbitrary nonlinear map.

The literature assessment checked on 21 August 2026 reports partial results for multiple-eigenvalue cases, including special Fu\v cik-branch analyses, while the requested general nonlinear spectral description remains unresolved in that assessment. It supplies no general class solving the second question. These are dated literature-triage claims requiring independent verification.
\subsection{Short English statement}
Describe what happens near a multiple nonlinear eigenvalue, and find structural conditions on homogeneous finite-dimensional maps that make every eigenvalue isolated. The source restores both questions but omits the operator and class hypotheses needed for a sharper theorem.
\subsection{Sources}
\begin{itemize}
\item \emph{Mini-Workshop: Nonlinear Spectral and Eigenvalue Theory} (2004), Oberwolfach Reports, DOI 10.4171/owr/2004/08, pp. 407--438, \url{https://doi.org/10.4171/owr/2004/08}.
\item Literature assessment checked 21 August 2026: ``Oberwolfach report,'' \url{https://oa.tib.eu/renate/server/api/core/bitstreams/aaddbeeb-c549-49a0-995d-8691ab5156bd/content}; ``Fuccik-spectrum paper,'' \url{https://arxiv.org/abs/1407.0190}; ``Recent discrete Fuccik paper,'' \url{https://arxiv.org/abs/2412.11709}.
\item Local UnsolvedMath registry, permanent ID $30000032$, exact source code \texttt{catalogue entry}, created 25 August 2026 and updated 16 September 2026. The clean statement is labelled corrected and verified; supplied verification explicitly restores both questions without inventing the operator.
\end{itemize}
\end{document}
